\chapter{From Starting Finance Club PoliTo Research to the Integrated Barrick Project}
\label{ch:introduction}

\section{Research question, community and valuation object}

% Sources: CONTENT-LEDGER BG-U-01, BG-U-10--14; CL-001/002/011/013;
% Team 8 T8-17--19; Team 4 MAP-T4-08--12; Team 5 T5-02--19;
% Research-area guide, project organisation and unification workflow.
Starting Finance Club PoliTo Research did not begin as a preassembled
quantitative-finance laboratory.
It originated as a comparatively vertical group, strongly represented by
management-engineering students and initially centred on events and financial
education. Research activity first appeared through short studies developed by
small groups over roughly monthly horizons. As contributors with different
academic paths entered the area, those isolated exercises gradually revealed a
larger opportunity: to build a common annual project in which every member
could contribute while learning how separate disciplines meet in a real
valuation problem.

That evolution matters because the project's main advantage is
interdisciplinarity. A mathematics curriculum can reasonably treat probability,
analysis or numerical foundations as prerequisites; a management curriculum
can do the same with accounting and valuation; a computing curriculum can
assume familiarity with implementation. The team deliberately avoided those
silent assumptions. The project was designed from the bottom up so that it
could extend, rather than merely repeat, the contributors' university
programmes and introduce models and techniques that many of them had not met
in formal courses.

The resulting chain joins economics and corporate finance with mathematics,
econometrics, statistics and probability, stochastic methods, numerical
methods, optimisation, operations research and computing. At university these
subjects are often taught as distinct modules. In the Barrick project they are
forced to interact: probability supplies the language of uncertainty;
econometrics tests the historical series; stochastic calculus defines the
continuous-time models; numerical methods make calibration and simulation
possible; optimisation estimates parameters; operational modelling converts
gold, production and costs into margins; and corporate finance closes the DCF.
The thesis is therefore both a valuation study and a record of how theoretical
ideas become mutually dependent when applied to one object.

The integrated research question is: \emph{how much does the choice of
gold-price model matter for a discounted operating-margin proxy of Barrick when
production, costs and discount-rate shocks are held fixed, and how can
option-implied information enter that valuation without being confused with a
physical operating forecast?} The object is not a point price produced by an
isolated option model. It is the conditional distribution of a discounted
operating-margin proxy for an eight-mine gold perimeter. A reconciled
enterprise value, and from it equity value and value per diluted share,
require separately verified accounting and capital bridges. Because these
bridges are not yet reconciled, this edition
reports aggregate operating proxies only, with no per-share value or
market-price comparison, and no target price or recommendation; the bridge and
any market comparison are future work.

This project was developed within Research, the Starting Finance Club PoliTo
area dedicated to collaborative quantitative, economic, financial and
methodological projects. Its members bring heterogeneous backgrounds---from
management engineering and mathematics to data science, computational methods
and quantitative finance---so the research process is organised around
knowledge sharing rather than an assumed common starting point.

That collaborative setting is material to the valuation design. The commodity
block must state whether paths are risk-neutral, physical or explicit
scenarios; the operating block must state production, cost and realised-price
units; the accounting block must reconcile operating margin, depreciation,
taxes and reinvestment; and the valuation block must state discounting,
terminal closure and every enterprise-to-equity adjustment. A result becomes
economically interpretable only when contributors working on different modules
use the same date, perimeter, currency, notation and manifest. The research
question is therefore also an integration question: whether heterogeneous
analytical work can be connected without hiding disagreements in definitions
or silently filling missing data.

\section{From eight autonomous studies to one chain of evidence}

% Sources: CONTENT-LEDGER authority table; MAP-T1--T4; T5-01--T8-20;
% Research-area guide, unification and review workflow.
The eight source studies are complementary modules
(Table~\ref{tab:team-question-model-output}), with different questions,
conventions and validation states, not chapters of one paper. Unification was
therefore a separate research phase rather than copy-and-paste assembly. Each
concept received one source authority; outputs were classified as
theoretical, synthetic, legacy empirical, current empirical or newly
generated; shared data and notation contracts fixed dates, currencies, units,
measures and interfaces; and source-model implementations were parity-tested
before being connected to the unified runner. The LSE (London Strategic Edge)
data refresh was kept as a separate evidence layer, so that a new market
observation could not be mistaken for a re-estimation of every legacy study.
Section~\ref{sec:team-study-synthesis} compares the resulting roles.

\section{Contribution, falsifiable boundaries and roadmap}

% Sources: VAL-001 gate sequence in GAP-REGISTER; CL-001/005/012/013.
The first contribution is an auditable handoff from option-implied dynamics to
corporate quantities. The second is an evidence hierarchy that separates
source-model parity, market-data validity, statistical fit, predictive tests,
numerical convergence and accounting reconciliation. The third is a reporting
contract: a stochastic valuation must expose its median and tails, its
numerical precision and the assumptions that move the distribution, rather
than presenting a single estimate with false precision.
Together, these contributions make disagreement inspectable: a reader can see
whether a result changes because of market evidence, model form, operating
assumptions, accounting treatment or the equity bridge.

Methodologically, Black--Scholes/GBM, Heston, Bates--Poisson and Full
Bates--Hawkes are calibrated, with nested refinements and the trust's 0.40\%
carry, to 64 GLD calls of 2 September 2026, and drive daily gold paths from the
XAUUSD spot on a dated clock that values only flows after that date. With
production, costs and WACC shocks common to all four branches, the median
proxies lie between USD~\ValMedianLow{} and \ValMedianHigh~billion, within the
numerical range of a single model (USD~\NumericalMedianLow--\NumericalMedianHigh~billion),
and the terminal value carries about \TerminalShare\% of the mean. In
\OosDates{} rolling repricing tests Full Bates--Hawkes has the lowest error but
no model beats implied-volatility persistence. The main limits are equally
explicit: the price laws are risk-neutral without a validated physical map,
operating uncertainty is not propagated, and the accounting and equity bridges
remain open.

The same structure creates falsifiable boundaries. Better option fit does not
prove better out-of-sample forecasts; better forecasts do not prove a valid
GLD-to-realised-gold bridge; a valid commodity bridge does not prove the mine
or accounting perimeter; and a coherent enterprise value does not determine
equity value until net debt, minorities, non-operating assets and diluted
shares are reconciled. Each boundary has a corresponding validation check in
Chapter~\ref{ch:unified-experiment}. Failure at a later check does not erase an earlier technical
result, but it narrows the conclusion that the integrated study may draw from
that result.

The thesis follows the same chain. Part I reconstructs the genesis of the
research, the monetary and macroeconomic role of gold, and the Barrick
valuation perimeter. Part II consolidates the mathematical, econometric and
numerical prerequisites from the bottom up. Part III first tests the legacy
gold series with discrete-time methods, then introduces the affine
stochastic-volatility and jump model ladder, and finally calibrates and
compares the models on the option surface. Part IV carries the gold paths and
the fixed production and cost forecasts through component margin and DCF,
validates the integrated experiment and closes with limitations and future
research. No part of the
study provides a buy, sell or hold recommendation.
